\section{Les deux grandes catégories de modèles et le diagramme de Gantt}
En Merise, on distingue deux grandes catégories de modèles : les modèles de traitement (MCC, MCT et MOT) et les modèles de données (MCD et MLD). Dans notre méthodologie de conception, nous avons étudié les trois différents modèles de traitements.
La représentation graphique des modèles est presentée a la fin de chaque sous section concernée et le diagramme de Gantt sera présenté dans l’Annexe~\ref{ann:diagrammes}.

\subsection{Modèle Conceptuel de Communication (MCC)}
Le MCC représente les acteurs impliqués dans le processus étudié ainsi que les flux d’informations échangés entre eux. 
Dans cette étude, il met en évidence les relations entre les acteurs internes et externes, notamment \textbf{le technicien, l’utilisateur, les fournisseurs, le bailleur et le capteur 8-en-1}. 

\begin{longtable}{|p{1cm}|p{6.5cm}|p{6.5cm}|}
    \caption{Représentation des flux}
    \label{tab:flux} \\
    \hline
    \textbf{N°} & \textbf{Émetteur $\rightarrow$ Récepteur} & \textbf{Libellé du flux} \\
    \hline
    \endfirsthead
    \hline
    \textbf{N°} & \textbf{Émetteur $\rightarrow$ Récepteur} & \textbf{Libellé du flux} \\
    \hline
    \endhead
    1 & Technicien $\rightarrow$ Fournisseur capteur & Commande de capteurs \\
    \hline
    2 & Fournisseur capteur $\rightarrow$ Technicien & Établissement et transmission du devis \\
    \hline
    3 & Technicien $\rightarrow$ Fournisseur capteur & Devis accepté \\
    \hline
    4 & Technicien $\rightarrow$ Bailleur & Demande de fonds (capteur) \\
    \hline
    5 & Bailleur $\rightarrow$ Technicien & Fonds mis à disposition (capteur) \\
    \hline
    6 & Technicien $\rightarrow$ Fournisseur capteur & Paiement \\
    \hline
    7 & Fournisseur capteur $\rightarrow$ Technicien & Livraison \\
    \hline
    8 & Technicien $\rightarrow$ Fournisseur intrant & Commande des intrants \\
    \hline
    9 & Fournisseur intrant $\rightarrow$ Technicien & Établissement et transmission du devis \\
    \hline
    10 & Technicien $\rightarrow$ Fournisseur intrant & Devis accepté \\
    \hline
    11 & Technicien $\rightarrow$ Bailleur & Demande de fonds (intrants) \\
    \hline
    12 & Bailleur $\rightarrow$ Technicien & Fonds mis à disposition (intrants) \\
    \hline
    13 & Technicien $\rightarrow$ Fournisseur intrant & Paiement \\
    \hline
    14 & Fournisseur intrant $\rightarrow$ Technicien & Livraison \\
    \hline
    15 & Technicien $\rightarrow$ Capteur 8-en-1 & Configuration du capteur \\
    \hline
    16 & Capteur 8-en-1 $\rightarrow$ Technicien & Envoi des données du sol \\
    \hline
    17 & Technicien $\rightarrow$ OpenWeather & Requête \\
    \hline
    18 & OpenWeather $\rightarrow$ Technicien & Envoi des données météorologiques \\
    \hline
    19 & Technicien $\rightarrow$ Utilisateur & Transmission des résultats et recommandations \\
    \hline
    20 & Utilisateur $\rightarrow$ Technicien & Retour \\
    \hline
\end{longtable}

\begin{figure}[H]
    \centering
    \includegraphics[width=\textwidth]{images/MCC.pdf}
    \caption{Modèle conceptuel de communication (MCC). }
    \label{fig:mcc}
\end{figure}

\subsection{Modèle Conceptuel de Traitement (MCT)} 
Le MCT décrit l’enchaînement des opérations du processus, depuis l’identification des besoins en équipements jusqu’à la transmission des résultats. 
Il met en évidence les \textbf{événements déclencheurs, les opérations réalisées, les conditions de leur exécution et les résultats obtenus}. 
Le tableau récapitulatif présente les principales opérations du processus.

{\small
\begin{longtable}{|p{0.8cm}|p{4.2cm}|p{4.2cm}|p{4.8cm}|}
    \caption{Correspondance entre les opérations, les événements déclencheurs et les résultats du MCT}
    \label{tab:operations_mct} \\
    \hline
    \textbf{N°} & \textbf{Opération} & \textbf{Événement(s) déclencheur(s)} & \textbf{Résultat} \\
    \hline
    \endfirsthead
    \hline
    \textbf{N°} & \textbf{Opération} & \textbf{Événement(s) déclencheur(s)} & \textbf{Résultat} \\
    \hline
    \endhead
    1 & Commander le capteur 8-en-1 auprès du fournisseur capteur & Besoin du capteur identifié et accord de l'utilisateur & En attente du devis \\
    \hline
    2 & Traitement du devis & Devis reçu & Acceptation ou refus ; si refus, fin de prestation \\
    \hline
    3 & Paiement & Acceptation du devis, demande de fonds et réception des fonds & Paiement effectué \\
    \hline
    4 & Installer et configurer le capteur & Capteur livré (ou reçu) & Capteur opérationnel \\
    \hline
    5 & Réaliser l'étude du sol (collecte des données du sol à l'aide du 8-en-1) & Accord de l'utilisateur & Collecte de données effectuée \\
    \hline
    6 & Analyse des résultats & Collecte effectuée & Besoin de correction identifié (si résultats valides) et ne pas se prononcer (si résultats invalides) \\
    \hline
    7 & Commande d'intrants & Accord de l'utilisateur & Attente du devis \\
    \hline
    8 & Traitement du devis & Devis reçu & Fin de prestation si devis refusé, ou demande de fonds si devis accepté \\
    \hline
    9 & Paiement & Fonds reçus & Paiement effectué \\
    \hline
    10 & Réceptionner les intrants & Livraison d'intrants & Intrants reçus \\
    \hline
    11 & Réaliser l'amendement sur la parcelle & Besoin de correction identifié et intrants reçus & Parcelle amendée \\
    \hline
    12 & Configuration et relevés de données (à chaque connexion du capteur) & Parcelle amendée & Données du sol reçues \\
    \hline
    13 & Requête AccuWeather & Parcelle amendée & Données météo reçues si connexion disponible, ou en attente de connexion si connexion non disponible \\
    \hline
    14 & Calculer les indicateurs agronomiques & Données du sol reçues et données météo reçues & Indicateurs agronomiques calculés \\
    \hline
    15 & Évaluer le besoin d'irrigation & Indicateurs agronomiques calculés & Irrigation nécessaire si $S_a < S_c$, ou stock suffisant (ne pas irriguer) si $S_a \geq S_c$ \\
    \hline
    16 & Irriguer la parcelle & Irrigation nécessaire et eau disponible & Irrigation réalisée \\
    \hline
    17 & Évaluer le besoin de fertilisation & Indicateurs agronomiques calculés & Fertilisation nécessaire si déficit de N/P/K, ou stock suffisant si pas de déficit \\
    \hline
    18 & Fertiliser la parcelle & Fertilisation nécessaire et disponibilité des intrants & Fertilisation réalisée \\
    \hline
    19 & Consolider les résultats et rédiger le rapport d'avancement & Irrigation réalisée, stock suffisant (ne pas irriguer), fertilisation réalisée et stock suffisant (ne pas fertiliser) & Rapport d'avancement rédigé \\
    \hline
    20 & Transmission du rapport & Rapport d'avancement rédigé & \\
    \hline
\end{longtable}
}

\begin{figure}[H]
    \centering
    \includegraphics[width=\textwidth, height=0.85\textheight, keepaspectratio]{images/MCT1.pdf}
    \caption{Modèle conceptuel de traitement (MCT), partie 1/6. }
    \label{fig:mct1}
\end{figure}

\begin{figure}[H]
    \centering
    \includegraphics[width=\textwidth, height=0.85\textheight, keepaspectratio]{images/MCT2.pdf}
    \caption{Modèle conceptuel de traitement (MCT), partie 2/6. }
    \label{fig:mct2}
\end{figure}

\begin{figure}[H]
    \centering
    \includegraphics[width=\textwidth, height=0.85\textheight, keepaspectratio]{images/MCT3.pdf}
    \caption{Modèle conceptuel de traitement (MCT), partie 3/6. }
    \label{fig:mct3}
\end{figure}

\begin{figure}[H]
    \centering
    \includegraphics[width=\textwidth, height=0.85\textheight, keepaspectratio]{images/MCT4.pdf}
    \caption{Modèle conceptuel de traitement (MCT), partie 4/6. }
    \label{fig:mct4}
\end{figure}

\begin{figure}[H]
    \centering
    \includegraphics[width=\textwidth, height=0.85\textheight, keepaspectratio]{images/MCT5.pdf}
    \caption{Modèle conceptuel de traitement (MCT), partie 5/6. }
    \label{fig:mct5}
\end{figure}

\begin{figure}[H]
    \centering
    \includegraphics[width=\textwidth, height=0.85\textheight, keepaspectratio]{images/MCT6.pdf}
    \caption{Modèle conceptuel de traitement (MCT), partie 6/6. }
    \label{fig:mct6}
\end{figure}

\subsection{Modèle Organisationnel de Traitement (MOT)}
Le MOT complète le MCT en précisant, pour chaque opération, les modalités organisationnelles de son exécution, notamment \textbf{le lieu, le moment et la nature} du traitement. 
Il permet ainsi de distinguer les opérations manuelles de celles envisagées comme automatisables.

\begin{figure}[H]
    \centering
    \includegraphics[width=\textwidth, height=0.85\textheight, keepaspectratio]{images/MOT1_1.pdf}
    \caption{Modèle organisationnel de traitement (MOT), partie 1/9. }
    \label{fig:mot1.1}
\end{figure}

\begin{figure}[H]
    \centering
    \includegraphics[width=\textwidth, height=0.85\textheight, keepaspectratio]{images/MOT1_2.pdf}
    \caption{Modèle organisationnel de traitement (MOT), partie 2/9. }
    \label{fig:mot1.2}
\end{figure}

\begin{figure}[H]
    \centering
    \includegraphics[width=\textwidth, height=0.85\textheight, keepaspectratio]{images/MOT1_3.pdf}
    \caption{Modèle organisationnel de traitement (MOT), partie 3/9. }
    \label{fig:mot1.3}
\end{figure}

\begin{figure}[H]
    \centering
    \includegraphics[width=\textwidth, height=0.85\textheight, keepaspectratio]{images/MOT1_4.pdf}
    \caption{Modèle organisationnel de traitement (MOT), partie 4/9. }
    \label{fig:mot1.4}
\end{figure}

\begin{figure}[H]
    \centering
    \includegraphics[width=\textwidth, height=0.85\textheight, keepaspectratio]{images/MOT1_5.pdf}
    \caption{Modèle organisationnel de traitement (MOT), partie 5/9. }
    \label{fig:mot1.5}
\end{figure}

\begin{figure}[H]
    \centering
    \includegraphics[width=\textwidth, height=0.85\textheight, keepaspectratio]{images/MOT1_6.pdf}
    \caption{Modèle organisationnel de traitement (MOT), partie 6/9. }
    \label{fig:mot1.6}
\end{figure}

\begin{figure}[H]
    \centering
    \includegraphics[width=\textwidth, height=0.85\textheight, keepaspectratio]{images/MOT1_7.pdf}
    \caption{Modèle organisationnel de traitement (MOT), partie 7/9. }
    \label{fig:mot1.7}
\end{figure}

\begin{figure}[H]
    \centering
    \includegraphics[width=\textwidth, height=0.85\textheight, keepaspectratio]{images/MOT1_8.pdf}
    \caption{Modèle organisationnel de traitement (MOT), partie 8/9. }
    \label{fig:mot1.8}
\end{figure}

\begin{figure}[H]
    \centering
    \includegraphics[width=\textwidth, height=0.85\textheight, keepaspectratio]{images/MOT1_9.pdf}
    \caption{Modèle organisationnel de traitement (MOT), partie 9/9. }
    \label{fig:mot1.9}
\end{figure}